\PassOptionsToPackage{numbers}{natbib}
\documentclass{article}
\usepackage{iclr2024_conference}
\usepackage{luatexja-fontspec}
\usepackage{hyperref}
\usepackage{url}
\usepackage{booktabs}
\usepackage{amsmath}
\usepackage{amsfonts}
\usepackage{graphicx}
\graphicspath{{../}}

\InputIfFileExists{values.tex}{}{}
\providecommand{\airasval}[1]{\textbf{??airasval:\detokenize{#1}??}}
\providecommand{\unverified}[1]{#1}
\providecommand{\airasrecordlink}[1]{#1}

\title{SciGym Table 1 の事前登録再現: 現行 API の Gemini 2.5 と GPT-4.1 で論文の平均 STE は戻るか}
\author{AIRAS}

\begin{document}
\maketitle

\begin{abstract}
SciGym \citep{duan-2025-measuring} は、反応をすべて取り除いた SBML モデルを LLM エージェントに渡し、摂動実験とコード実行を繰り返して反応機構を復元させるベンチマークである。論文の Table 1 は 6 つのフロンティアモデルを SciGym-small 137 件で各 1 回走らせた平均値を報告しているが、実行間の分散も、別環境・別時点での再現性も示していない。本研究は、公開コードの Controller と評価器をそのまま用い、現在 API で提供されている Gemini-2.5-Pro、Gemini-2.5-Flash、GPT-4.1、GPT-4.1-mini の 4 モデルで Table 1 の対応する 4 行を再現する。実験前に、各モデルの平均 Simulation Trajectory Error（STE）が論文値を 0.05 より大きくは上回らないことを反証線として登録し、予測区間を $\pm 0.05$ と宣言する。Reaction Matching Score（RMS）は同じ実行から表として報告する。実行は RIKYU（aarch64）上で行い、論文設定からの逸脱（Gemini の GA 版、libroadrunner 2.7.0、応答トークン上限）はすべて実験前に宣言する。
\end{abstract}

\section{はじめに}
\label{sec:intro}
LLM に科学的発見のサイクル全体（実験の設計、結果の解析、機構の提案）を行わせて測るベンチマークとして、SciGym \citep{duan-2025-measuring} が提案された。BioModels 由来の SBML モデルから反応をすべて取り除いたものをエージェントに渡し、エージェントは初期濃度を変えた摂動実験を要求し、返ってきた時系列を Python で解析し、最終的な SBML を提出する。提出モデルは真のモデルと、軌道誤差（STE）と反応の一致度（RMS）で比較される。

論文の主結果である Table 1 は、3 社の pro / mini 計 6 モデルを SciGym-small の 137 件で評価した平均値である \cite[s1.p10]{duan-2025-measuring}。この評価は各モデル 1 回の実行であり、temperature 1.0 でサンプリングされる \cite[s2.p1]{h4duan-2025-scigym}。したがって Table 1 の数値がどの程度安定なのか、また同じコードを別の環境と時点で走らせて同じ数値が戻るのかは、論文からは分からない。さらに Gemini の 2 モデルはプレビュー版で評価されており \cite[s1.p1]{duan-2025-measuring}、現在は同じ版を呼び出せない。

本研究はこのギャップに対し、Table 1 のうち Gemini と OpenAI の 4 行を、公開コード \citep{h4duan-2025-scigym} の Controller をそのまま使って再現する。Claude の 2 モデルは本研究が用いる API ゲートウェイで提供されていないため対象外とする。判定は AIRAS の事前登録の枠組みに従い、仮説・判定基準・予測区間・実行条件を実験前に記録へ凍結し、実験後は計測値だけを流し込む。

\section{関連研究}
\label{sec:related}
SciGym \citep{duan-2025-measuring} は、単一スキルを測る既存のベンチマークと異なり、エージェントが自ら摂動実験を設計してその結果を受け取る閉ループを評価する。評価指標は 3 つで、種間の有向関係の F1（Network Topology Score）、反応物と生成物の集合が一致した反応の適合率・再現率・F1（RMS。modifier まで要求する厳格版もある）\cite[s1.p8]{duan-2025-measuring}、および全種の時系列の対称平均絶対パーセント誤差（STE）\cite[s1.p9]{duan-2025-measuring} である。本研究は公開コード \citep{h4duan-2025-scigym} が実行時に書き出す STE と RMS を対象とし、Network Topology Score は実行時に計算されないため扱わない。

論文は pro モデルが mini モデルを一貫して上回り \cite[s1.p6]{duan-2025-measuring}、Gemini-Pro が最良であった \cite[s1.p7]{duan-2025-measuring} と報告している。本研究はこれらの順序関係についても、判定基準には含めないが、得られた表から読み取って報告する。

\section{仮説と予測}
\label{sec:hypothesis}
以下は記録（record.json）から生成された仮説・主張・判定基準・予測区間の一覧である。実験後は Observed と Verdict が計測値から埋められる。

% AUTO-GENERATED by the update_record tool. DO NOT EDIT.
% verify_paper regenerates this file from record.json and the run outputs
% and fails on any difference, so manual edits always surface.
\noindent\textbf{Hypothesis 1.} SciGym 論文の Table 1 のうち Gemini-2.5-Pro、Gemini-2.5-Flash、GPT-4.1、GPT-4.1-mini の 4 行は、公開コードの Controller をそのまま用い、現在 API で提供されている同名モデル（Gemini 2 種は GA 版、GPT-4.1 2 種は論文と同じ 2025-04-14 スナップショット）で SciGym-small 137 件を最大 20 反復で解かせたとき、各モデルの平均 Simulation Trajectory Error（STE）が論文値から 0.05 以内に収まる形で再現される。~\cite[s1.p1, s1.p5, s1.p10]{duan-2025-measuring}
\begin{enumerate}
\item[\textbf{Claim 1}] 現行の Gemini-2.5-Pro（google/gemini-2.5-pro）で SciGym-small 137 件を最大 20 反復で解かせたときの平均 STE は、論文の 0.3212 を 0.05 より大きくは上回らない。~\cite[s1.p1, s1.p2, s1.p3, s1.p4, s1.p5, s1.p7, s1.p10]{duan-2025-measuring}, \cite[s2.p1]{h4duan-2025-scigym}
  \emph{Rationale:} Table 1 で最良とされた行がそのまま再現されることは、論文の主結果が特定のプレビュー版や実行環境に依存していないことの直接の証拠であり、h1 の Gemini-2.5-Pro 行を担う。
  \emph{Criterion:} \texttt{\detokenize{gemini-2.5-pro}}.\texttt{\detokenize{ste}} $-$ 0.3212 $\leq 0.05$.
  \emph{Prediction:} $[-0.05, 0.05]$ (論文 Table 1 の 1 回実行の平均値。137 件平均の標準誤差はおよそ 0.02 と見積もり、モデル版の差を含めても ±0.05 に収まると予測する).
  \emph{Observed:} 0.0935808.
  \emph{Verdict:} refuted.
\item[\textbf{Claim 2}] 現行の Gemini-2.5-Flash（google/gemini-2.5-flash）で SciGym-small 137 件を最大 20 反復で解かせたときの平均 STE は、論文の 0.4181 を 0.05 より大きくは上回らない。~\cite[s1.p1, s1.p2, s1.p3, s1.p4, s1.p5, s1.p6, s1.p10]{duan-2025-measuring}, \cite[s2.p1]{h4duan-2025-scigym}
  \emph{Rationale:} Gemini 系の mini 変種の行が再現されることで、h1 の Gemini-2.5-Flash 行を担い、c1 と合わせて Gemini 系の pro / mini の関係が論文どおりかを読めるようにする。
  \emph{Criterion:} \texttt{\detokenize{gemini-2.5-flash}}.\texttt{\detokenize{ste}} $-$ 0.4181 $\leq 0.05$.
  \emph{Prediction:} $[-0.05, 0.05]$ (論文 Table 1 の 1 回実行の平均値。137 件平均の標準誤差はおよそ 0.02 と見積もり、モデル版の差を含めても ±0.05 に収まると予測する).
  \emph{Observed:} 0.0284183.
  \emph{Verdict:} supported.
\item[\textbf{Claim 3}] GPT-4.1（openai/gpt-4.1、2025-04-14 スナップショット）で SciGym-small 137 件を最大 20 反復で解かせたときの平均 STE は、論文の 0.4611 を 0.05 より大きくは上回らない。~\cite[s1.p1, s1.p2, s1.p3, s1.p4, s1.p5, s1.p10]{duan-2025-measuring}, \cite[s2.p1]{h4duan-2025-scigym}
  \emph{Rationale:} 論文と同じ日付スナップショットが使える唯一の pro モデルであり、モデル版の差を含まない最も純粋な再現条件として h1 の GPT-4.1 行を担う。
  \emph{Criterion:} \texttt{\detokenize{gpt-4.1}}.\texttt{\detokenize{ste}} $-$ 0.4611 $\leq 0.05$.
  \emph{Prediction:} $[-0.05, 0.05]$ (論文 Table 1 の 1 回実行の平均値。同じスナップショットなので差は実行の揺らぎとシミュレータ差のみで、137 件平均の標準誤差およそ 0.02 の範囲に収まると予測する).
  \emph{Observed:} 0.0469407.
  \emph{Verdict:} supported.
\item[\textbf{Claim 4}] GPT-4.1-mini（openai/gpt-4.1-mini、2025-04-14 スナップショット）で SciGym-small 137 件を最大 20 反復で解かせたときの平均 STE は、論文の 0.6007 を 0.05 より大きくは上回らない。~\cite[s1.p1, s1.p2, s1.p3, s1.p4, s1.p5, s1.p6, s1.p10]{duan-2025-measuring}, \cite[s2.p1]{h4duan-2025-scigym}
  \emph{Rationale:} 同じスナップショットが使える mini モデルの行を担い、c3 と合わせて OpenAI 系の pro / mini の関係が論文どおりかを読めるようにする。
  \emph{Criterion:} \texttt{\detokenize{gpt-4.1-mini}}.\texttt{\detokenize{ste}} $-$ 0.6007 $\leq 0.05$.
  \emph{Prediction:} $[-0.05, 0.05]$ (論文 Table 1 の 1 回実行の平均値。同じスナップショットなので差は実行の揺らぎとシミュレータ差のみで、137 件平均の標準誤差およそ 0.02 の範囲に収まると予測する).
  \emph{Observed:} 0.00101829.
  \emph{Verdict:} supported.
\end{enumerate}
\noindent\emph{Assumed, so that the claims together imply Hypothesis 1:}
\begin{itemize}
\item 同名モデルの現行版（Gemini-2.5-Pro / Flash の GA 版、GPT-4.1 / 4.1-mini の 2025-04-14 スナップショット）が、論文で使われた版と同等の能力を持つこと（c1〜c4。Gemini はプレビュー版が提供終了のため GA 版で代替する）
\item aarch64 向けの libroadrunner 2.7.0 と antimony 2.14.0 が、論文環境の libroadrunner 2.8.0 と同じ軌道と評価値を与えること（c1〜c4）
\item temperature 1.0 での 1 回実行の平均 STE（137 件平均）の実行間変動が 0.05 より十分小さいこと（c1〜c4）
\item Vercel AI Gateway の OpenAI 互換エンドポイント経由の各モデルが、各社の公式 API と同じ応答分布を持つこと（c1〜c4）
\item 平均 STE が Table 1 の再現性を代表すること。RMS（modifier あり / なし）は同じ実行から表として報告するが、判定基準には用いない（c1〜c4）
\end{itemize}


\paragraph{判定基準と予測区間の根拠.}
各主張の反証線は「（再現の平均 STE）$-$（論文値）$\le 0.05$」である。0.05 は、論文 Table 1 で隣り合うモデルの最小差（Gemini-2.5-Pro の \unverified{0.3212} と Claude-3.7-Sonnet の \unverified{0.3615} の差 \unverified{0.04}）とほぼ同じで、これより大きく外れると Table 1 の順位付けの粒度で再現に失敗したことになる。1 件の STE は $[0, 1]$ の値をとり、標準偏差を 0.25 程度と見積もると 137 件平均の標準誤差は約 0.02、2 つの 1 回実行の差の標準誤差は約 0.03 である。予測区間は $\pm 0.05$ とする。基準は悪化方向の片側なので、現行版のモデルが論文より良い結果を出しても主張は成立するが、予測区間を下に外れた場合はそのとおり報告する。

RMS の適合率・再現率・F1 も同じ実行から表として報告し、各モデルの F1 が論文値 $\pm 0.05$ に入ることを予測するが、1 つの主張に置ける判定基準は 1 つであるため、判定には用いない。

\section{方法}
\label{sec:method}
\paragraph{タスク.}
各インスタンスは真の SBML（\texttt{truth.xml}）、反応を取り除いた不完全な SBML（\texttt{partial.xml}）、シミュレーション条件（\texttt{truth.sedml}）からなる。エージェントは不完全な SBML を受け取り、各ターンで実験（初期濃度の変更を指定し、真のモデルの時系列を得る）、コード実行（numpy / pandas / scipy / libsbml と、仮説モデルを走らせる \texttt{simulate} 関数が使える）、提出のいずれかを markdown 形式で返す。提出は任意の時点で行え、反復数の上限は 20 である \cite[s1.p3]{duan-2025-measuring}。提出した SBML が無効な場合は追加のデバッグ反復が許され、それでも無効なら不完全 SBML のまま採点される \cite[s1.p4]{duan-2025-measuring}。

\paragraph{実行系.}
ループ、プロンプト、実験の実行、提出の採点はすべて公開コード \citep{h4duan-2025-scigym} の \texttt{Controller} と \texttt{Evaluator} をそのまま用いる。本研究が書くのは、(1) OpenAI 互換 API でモデルを呼ぶ \texttt{LLM} 実装（公式の GPT 実装と同じく、システムプロンプトと全履歴を毎回送る）、(2) インスタンスごとに 1 プロセスで Controller を起動し、137 件を並列に走らせる駆動部、(3) 各インスタンスの \texttt{evaluation.json} を集計する部分だけである。公式コードでは 5 つの履歴がクラス属性として定義されており 1 プロセスで複数件を走らせると履歴が持ち越されるが、1 件 1 プロセスとすることでこの問題を避ける。

\paragraph{指標.}
各インスタンスの \texttt{evaluation.json} の \texttt{observe\_smape} を STE、\texttt{rp\_precision} / \texttt{rp\_recall} / \texttt{rp\_f1} を modifier なしの RMS、\texttt{rpm\_*} を modifier ありの RMS とし、137 件の単純平均を各 run の \texttt{ste}、\texttt{rms\_*}、\texttt{rmsm\_*} とする。提出 SBML が無効で \texttt{success=false} となった件も、公式評価どおり不完全 SBML の採点値で平均に含める。補助的に、真の欠損反応と提出された追加反応の集合を 2 値分類として airas-eval に渡し、全件をプールした適合率・再現率・F1（micro）を照合用に記録する。

\section{実験設計}
\label{sec:design}
\paragraph{モデル.}
論文の 6 モデルのうち、Gemini-2.5-Pro-Preview-03-25、Gemini-2.5-Flash-Preview-04-17、GPT-4.1-2025-04-14、GPT-4.1-Mini-2025-04-14 \cite[s1.p1]{duan-2025-measuring} に対応する現行モデルとして、Vercel AI Gateway の OpenAI 互換エンドポイント経由で \texttt{google/gemini-2.5-pro}、\texttt{google/gemini-2.5-flash}、\texttt{openai/gpt-4.1}、\texttt{openai/gpt-4.1-mini} を用いる。GPT-4.1 の 2 つは論文と同じ 2025-04-14 スナップショットである。Gemini の 2 つはプレビュー版が提供終了しているため GA 版を用いる。これが本研究の再現条件の最大の違いであり、事前に宣言する。

\paragraph{データ.}
SciGym-small の 137 件 \cite[s1.p2]{duan-2025-measuring}（Hugging Face \texttt{h4duan/scigym-sbml} の small split）を用いる。

\paragraph{設定.}
\texttt{max\_iterations} は論文どおり 20（公開 config の既定値 5 \cite[s2.p2]{h4duan-2025-scigym} は動作確認用）、\texttt{eval\_debug\_rounds} は論文本文の 3 \cite[s1.p4]{duan-2025-measuring}（公開 config の既定値は 5 \cite[s2.p3]{h4duan-2025-scigym}）、\texttt{temperature} は公開 config の 1.0 \cite[s2.p1]{h4duan-2025-scigym}、摂動は初期濃度の変更のみとする。応答のトークン上限は 32768 とする。公式の GPT 実装の既定値 8192 では、思考トークンを出すモデルが本文を返せずに失敗しうるためである。各モデル、各インスタンスを 1 回ずつ実行する。

\paragraph{計算資源と環境.}
実行は Seyval を通じて RIKYU（Slurm、GB200 ノード、aarch64）で行う。GPU は使わないが、資源の単位が GPU であるため 1 GPU を要求する。環境は Dockerfile と \texttt{uv.lock} で固定する。aarch64 向けの wheel が無いため、libroadrunner は公式の 2.8.0 ではなく 2.7.0、antimony は 2.14.0 を用い、rrplugins と phrasedml は除外する。この差でシミュレーション結果がわずかに変わる可能性があり、仮説の前提として記録する。137 件は 1 件 1 プロセスで並列に実行する。

\paragraph{段階.}
sanity は 1 件（BIOMD0000000027）を 2 反復で走らせて配管を確かめる。pilot は sorted した 137 件から 14 件おきに選んだ 10 件を 20 反復で走らせ、1 件あたりの所要時間と API コストを測って全件実行の可否を判断する。主張の判定は full（137 件）の結果だけで行い、sanity と pilot の結果は記録に取り込まない。

\paragraph{run と主張の対応.}
\texttt{gemini-2.5-pro} $\to$ C1、\texttt{gemini-2.5-flash} $\to$ C2、\texttt{gpt-4.1} $\to$ C3、\texttt{gpt-4.1-mini} $\to$ C4。

% ===== airas: post-experiment sections, filled once runs exist =====

\section{結果}
\label{sec:results}
表~\ref{tab:table1} と図~\ref{fig:charts} に、4 モデルそれぞれの SciGym-small 137 件に対する平均 STE と RMS を示す。各値は run の \texttt{metrics.json} から記録経由で流し込んだものである。

% AUTO-GENERATED by the update_record tool. DO NOT EDIT.
% verify_paper regenerates this file from record.json and the run outputs
% and fails on any difference, so manual edits always surface.
\begin{table}[t]
\centering
\caption{SciGym-small 137 件における平均 STE（低いほど良い）と Reaction Matching Score（高いほど良い）。論文 Table 1 の Gemini / OpenAI の 4 行に対応する。}
\label{tab:table1}
\begin{tabular}{lrrrrrrr}
\hline
 & STE $\downarrow$ & RMS w/ mod. P & RMS w/ mod. R & RMS w/ mod. F1 & RMS w/o mod. P & RMS w/o mod. R & RMS w/o mod. F1 \\
\hline
Gemini-2.5-Flash & 0.4465 & 0.1599 & 0.1226 & 0.1344 & 0.2836 & 0.2327 & 0.2473 \\
GPT-4.1-mini & 0.6017 & 0.1839 & 0.1412 & 0.1511 & 0.2915 & 0.2551 & 0.2588 \\
Gemini-2.5-Pro & 0.4148 & 0.1707 & 0.1309 & 0.1426 & 0.3144 & 0.2754 & 0.2844 \\
GPT-4.1 & 0.5080 & 0.1925 & 0.1542 & 0.1614 & 0.3106 & 0.2670 & 0.2718 \\
\hline
\end{tabular}
\end{table}


\begin{figure}[t]
  \centering
  \includegraphics[width=0.48\linewidth]{images/chart/ste_by_model.png}
  \hfill
  \includegraphics[width=0.48\linewidth]{images/chart/rms_f1_by_model.png}
  \caption{4 モデルの平均 STE（左、低いほど良い）と平均 RMS F1（右、高いほど良い）。値は各 run の \texttt{metrics.json} から解決され、記録に宣言した仕様から再描画して検証される。}
  \label{fig:charts}
\end{figure}

\paragraph{主張の判定.}
各主張の判定基準は「平均 STE $-$ 論文値 $\le 0.05$」である（第\ref{sec:hypothesis}節）。論文 Table 1 の値 \cite[s1.p5]{duan-2025-measuring} と比べると、
GPT-4.1-mini の平均 STE は \airasval{gpt-4.1-mini.ste}（論文値 \unverified{0.6007}）、
GPT-4.1 は \airasval{gpt-4.1.ste}（論文値 \unverified{0.4611}）、
Gemini-2.5-Flash は \airasval{gemini-2.5-flash.ste}（論文値 \unverified{0.4181}）であり、
C4、C3、C2 はいずれも基準を満たし、差は予測区間 $[-0.05, 0.05]$ の内側に入った。ただし GPT-4.1 の差は約 \unverified{0.047} で、基準の 0.05 にほぼ張り付いている。
一方 Gemini-2.5-Pro の平均 STE は \airasval{gemini-2.5-pro.ste}（論文値 \unverified{0.3212}）で、差は約 \unverified{0.094} と基準を超え、C1 は反証された。

\paragraph{Gemini-2.5-Pro の実行で起きたこと.}
C1 の反証は、モデルの能力の差だけでは説明できない。gemini-2.5-pro の run では、実行途中で Vercel AI Gateway のチーム予算上限に達し、以降の API 呼び出しが HTTP 402 で拒否された。その結果、137 件のうち \airasval{gemini-2.5-pro.n_timed_out} 件が 3 回の試行とも途中で終わり、公式評価の「有効な SBML が提出されなかった」場合と同じく不完全 SBML の採点値（\texttt{success=false}）で平均に含まれている（\texttt{n\_success} は \airasval{gemini-2.5-pro.n_success}）。不完全 SBML の STE はモデルごとに 0.3〜0.6 程度になるため、この 25 件が平均を押し上げている。この run をそのまま結果として確定する判断は利用者が行い、記録の notes に追記した。参考として、完走した 112 件だけの平均 STE を \texttt{instances.json} から計算すると \unverified{0.3514} であり、これは論文値との差が 0.030 で基準の内側に入るが、事前登録の判定基準は 137 件平均で定義したため判定には用いない。他の 3 モデルは予算超過の前に完了しており、この影響を受けていない（gemini-2.5-flash では 1 件だけ提出 SBML が無効で \texttt{success=false}、\texttt{n\_success} は \airasval{gemini-2.5-flash.n_success}）。

\paragraph{RMS と順序関係.}
modifier なしの RMS F1 は、Gemini-2.5-Pro \airasval{gemini-2.5-pro.rms_f1}（論文値 \unverified{0.3383}）、GPT-4.1 \airasval{gpt-4.1.rms_f1}（\unverified{0.3038}）、GPT-4.1-mini \airasval{gpt-4.1-mini.rms_f1}（\unverified{0.2322}）、Gemini-2.5-Flash \airasval{gemini-2.5-flash.rms_f1}（\unverified{0.2005}）であった。pro 2 モデルは論文より低く、mini 2 モデルは論文より高い。差はいずれも 0.06 以内で、事前に予測した $\pm 0.05$ を pro 側でわずかに外れている。
順序関係については、平均 STE の順位（Gemini-2.5-Pro $<$ Gemini-2.5-Flash $<$ GPT-4.1 $<$ GPT-4.1-mini）と modifier なし RMS F1 の順位（Gemini-2.5-Pro $>$ GPT-4.1 $>$ GPT-4.1-mini $>$ Gemini-2.5-Flash）は論文 Table 1 と同じで、pro が mini を上回るという論文の主張 \cite[s1.p6]{duan-2025-measuring} と、Gemini-Pro が最良という主張 \cite[s1.p7]{duan-2025-measuring} は、Gemini-2.5-Pro の 25 件が不利に採点された状態でも保たれた。modifier ありの RMS F1 は GPT-4.1 \airasval{gpt-4.1.rmsm_f1} が最も高く、Gemini-2.5-Pro \airasval{gemini-2.5-pro.rmsm_f1} は 3 位で、この順位は論文と異なる。

\paragraph{完走件数とコスト.}
表~\ref{tab:tokens} に各 run の完走件数とトークン消費を示す。Gemini の出力トークンには思考トークンが含まれる。4 モデルの pilot と全件を合わせたゲートウェイの課金は約 \unverified{200} USD で、gemini-2.5-pro の全件実行の途中でチーム予算上限に達した。

% AUTO-GENERATED by the update_record tool. DO NOT EDIT.
% verify_paper regenerates this file from record.json and the run outputs
% and fails on any difference, so manual edits always surface.
\begin{table}[t]
\centering
\caption{各 run の完走件数とトークン消費。n\_success は有効な SBML が提出された件数（137 件中）。gemini-2.5-pro の不足分 25 件はゲートウェイの予算超過（HTTP 402）で 3 試行とも終わらず、不完全 SBML の採点値で代用した件。Gemini の出力トークンには思考トークンを含む。}
\label{tab:tokens}
\begin{tabular}{lrrr}
\hline
 & n\_success & 入力トークン & 出力トークン \\
\hline
Gemini-2.5-Pro & 112 & 32346906 & 3964227 \\
Gemini-2.5-Flash & 136 & 48858545 & 7738168 \\
GPT-4.1 & 137 & 38206200 & 1229341 \\
GPT-4.1-mini & 137 & 38253239 & 873529 \\
\hline
\end{tabular}
\end{table}


\section{考察}
\label{sec:discussion}
\paragraph{同一スナップショットの GPT-4.1 系は再現した.}
論文と同じ 2025-04-14 スナップショットを使えた GPT-4.1 と GPT-4.1-mini は、平均 STE が基準の内側に入った。GPT-4.1-mini の差は 0.001 とほぼ一致し、GPT-4.1 は 0.047 で基準ぎりぎりであった。1 回実行の平均 STE の変動を約 0.03 と見積もった事前の予測からすると、GPT-4.1 の差は偶然の範囲としては大きめであり、aarch64 向け libroadrunner 2.7.0 の差や、応答トークン上限の違い（公式 8192、本研究 32768）が影響した可能性は残る。

\paragraph{GA 版 Gemini は Flash では再現し、Pro は判定できなかった.}
Gemini-2.5-Flash はプレビュー版から GA 版に変わっても基準の内側に入り、RMS F1 は論文より高かった。Gemini-2.5-Pro は予算超過で 25 件が失敗したため、C1 の反証がモデル版の差によるものか、実行側の事故によるものかを本研究のデータでは分離できない。完走 112 件の参考値（STE \unverified{0.3514}）は基準の内側にあり、後者の寄与が大きいと考えられるが、これは事後の部分集合の平均であり、事前登録の判定を置き換えるものではない。予算を増やして 25 件を再実行すれば分離できる。記録は追記のみなので、再実行の結果は別の結果として同じ run に追加され、C1 の判定は変わらないまま併記される。

\paragraph{順序関係は保たれた.}
平均 STE と modifier なし RMS F1 の 4 モデルの順位は論文と一致した。Table 1 の主張のうち「pro が mini を上回る」「Gemini-Pro が最良」という質的な結論は、モデル版・シミュレータ版・実行環境が変わっても保たれたと言える。一方、絶対値は STE で 0.001〜0.094、RMS F1 で 0.03〜0.06 動いた。temperature 1.0 の 1 回実行である以上、Table 1 の値を 0.05 より細かい精度で比較することには意味がないというのが本研究の読みである。

\paragraph{実行上の教訓.}
LLM が書いたコードが \texttt{simulate()} で硬い ODE を解こうとすると C ライブラリ内で止まり、公式 Controller の 3 分制限（SIGALRM）が効かない場合があった。1 件 1 プロセスとプロセス単位の強制タイムアウトは必須である。また、外部 API の予算上限は run の途中で静かに結果を汚染する。402 のような課金エラーは他の失敗と区別して run を止める設計にすべきであった。

\section{結論}
\label{sec:conclusion}
SciGym Table 1 の Gemini / OpenAI の 4 行を、公開コードの Controller と現行 API のモデルで事前登録のうえ再現した。GPT-4.1、GPT-4.1-mini、Gemini-2.5-Flash の 3 モデルは平均 STE が論文値 $+0.05$ 以内に収まり、主張 C2〜C4 は支持された。Gemini-2.5-Pro は実行中の API 予算超過で 25 件が失敗したまま確定したため C1 は反証されたが、完走分の参考値は基準の内側にあり、4 モデルの順位は STE・RMS F1 とも論文と一致した。Table 1 の質的な結論は再現し、絶対値は 1 回実行の揺らぎと実行環境の差の範囲で動く、というのが本研究の結論である。

\section*{データの利用可能性}
記録は \airasrecordlink{record.json} にある。実験コード、環境の固定（Dockerfile と \texttt{uv.lock}）、各 run の結果と provenance は同じリポジトリに含まれる。

\bibliographystyle{iclr2024_conference}
\bibliography{references}

\end{document}
